\section{Arithmetic circuits with retained quantum reductions}
\label{sec:arithmetic-processes}

A field extension supplies two distinct quantum embeddings: inclusion of its
basis labels, and a normalized sum over each relative-trace fibre. Their
adjoints are successful reductions to the smaller field. To compose these
reductions as physical operations, we retain their failure outcomes and the
quantum system on which each outcome lives. Independent tensor composition
then preserves partial success on either system.

\begin{definition}[Named finite arithmetic registers]
Fix a prime $p$ and a named primitive root $\zeta_p\in\mathbb C$. For a finite
field $E/\mathbb F_p$, put
\[
 H_E=\ell^2(E),\qquad
 \psi_E(x)=\zeta_p^{\operatorname{Tr}_{E/\mathbb F_p}(x)},\qquad
 W_E(a,b)=Z_E(-b)X_E(a),
\]
where $X_E(a)|x\rangle=|x+a\rangle$ and
$Z_E(b)|x\rangle=\psi_E(bx)|x\rangle$. The attached prime-field symplectic
space is $E\oplus E$ with form
\[
 \Omega_E((a,b),(a',b'))=
 \operatorname{Tr}_{E/\mathbb F_p}(ab'-a'b).
\]
Choose a finite nonempty collection $\mathcal F$ of explicit finite fields
and a finite category $\mathcal I$ of named injective field maps between
them, containing identities and composites. Equality of field arrows means
equality of their functions on the named finite sets. The polynomial
presentations, element names and embeddings are retained choices. Use the
coefficient field $R_p=\mathbb Q(\zeta_p,\sqrt p)$, with its given complex
embedding, positive square root and complex conjugation.

For $i:K\hookrightarrow E$, let $T_i:E\to K$ be relative trace followed
by $i^{-1}$, and set
\[
 J_i|a\rangle=|i(a)\rangle,\qquad
 V_i|a\rangle=\frac{1}{\sqrt{|\ker T_i|}}
                  \sum_{T_i(x)=a}|x\rangle.
\]
The support code is $C_i=\operatorname{ran}J_i$, its projector is
$P_i=J_iJ_i^*$, and its phase stabilizer labels form
$N_i=\{0\}\oplus\ker T_i$. Its logical symplectic space is
$N_i^\perp/N_i$, identified with $K\oplus K$ by
$[(i(a),b)]\mapsto(a,T_i(b))$.

The quantum atoms are $Q_E$ and the ambient-labeled code atoms $C_i$.
A quantum word is a finite ordered sequence of atoms, with the empty word
$\mathbf1$ allowed; tensor is concatenation. Attach the symplectic space
$E\oplus E$ to $Q_E$, the quotient $N_i^\perp/N_i$ to $C_i$, and orthogonal
direct sums to words. An isotropic context is a flag of prime-field
isotropic subspaces in the attached space. Contexts remain geometric data,
without quotienting by a symmetry group; no generator for every context is
stipulated. A code atom retains its ambient field, embedding and isotropic
reduction label. No identification of a tensor of field atoms with one
extension-field atom is imposed.
\end{definition}
\begin{scope}
Characteristic two is included. There $R_p$ is real, so this is the stated
arithmetic gate fragment rather than the full Clifford category. All field
models and embeddings are named choices; they are not reconstructed from an
unmarked matrix algebra. Contexts are not the sole objects.
\end{scope}
\provenance{D1301,D1303--D1305,D1321}{definitions.md; theory/sidequests/frobenius-hierarchy/foundation.md}{theory/checks/frobenius\_hierarchy\_check.py (A1--A3)}{theory/verdicts/frobenius-hierarchy-category-r1.md}

\begin{definition}[Presented arithmetic amplitudes]
On these quantum words, take finite $R_p$-linear combinations of well-typed
circuits generated by identities, swaps, composition, tensor, dagger, and
the arrows
\[
\begin{array}{c|c}
 \text{arrow}&\text{source and target}\\ \hline
 w_E(a,b),f_E,u_E&Q_E\longrightarrow Q_E\\
 m_{E,d}&Q_E^{\otimes(d+1)}\longrightarrow Q_E^{\otimes(d+1)},\quad d\geq1\\
 j_i,v_i&Q_K\longrightarrow Q_E\\
 e_E&\mathbf1\longrightarrow Q_E\\
 t_i&Q_K\longrightarrow C_i\\
 c_i&C_i\longrightarrow Q_E .
\end{array}
\]
Dagger reverses types and conjugates coefficients. Coefficient scalars in
$R_p$ act centrally; general endomorphisms of the empty word can also contain
closed circuit diagrams and are not asserted to be coefficient scalars.
Equality is the least typed congruence generated by the linear dagger strict
symmetric monoidal axioms and the following equations and their daggers.
Write $r=[E:\mathbb F_p]$, and use $i:K\to L$, $j:L\to E$ for towers.
\begin{enumerate}
\item $w_E(0,0)=1$, each $w_E(a,b)$ is unitary, and
\[
 w_E(a,b)w_E(a',b')=\psi_E(ab')w_E(a+a',b+b').
\]
\item $f_E,u_E,m_{E,d}$ are unitary,
$f_E^4=1$, $u_E^r=1$, $m_{E,d}^{p}=1$, and $f_Eu_E=u_Ef_E$.
\item Frobenius satisfies
\[
 u_Ew_E(a,b)=w_E(a^p,b^p)u_E,\qquad
 u_E^{\otimes(d+1)}m_{E,d}=m_{E,d}u_E^{\otimes(d+1)}.
\]
\item $j_i^\dagger j_i=v_i^\dagger v_i=1$,
$j_{\mathrm{id}}=v_{\mathrm{id}}=1$,
$j_jj_i=j_{ji}$ and $v_jv_i=v_{ji}$.
\item For $i:K\to E$,
\[
 u_Ej_i=j_iu_K,\qquad u_Ev_i=v_iu_K,\qquad f_Ej_i=v_if_K.
\]
\item For $a\in K$ and $b\in E$,
\[
 w_E(i(a),b)j_i=j_iw_K(a,T_i(b)),\qquad
 m_{E,d}j_i^{\otimes(d+1)}=j_i^{\otimes(d+1)}m_{K,d}.
\]
\item $e_E^\dagger e_E=1$, $t_i^\dagger t_i=t_it_i^\dagger=1$,
and $c_it_i=j_i$.
\item With $b_{E,a}=w_E(a,0)e_E$,
\[
 b_{E,a}^\dagger b_{E,a'}=\delta_{a,a'},\qquad
 \sum_{a\in E}b_{E,a}b_{E,a}^\dagger=1.
\]
\end{enumerate}
The trace for the composite embedding is the transported composite trace.
There are no further arithmetic equations by default. Scalars are retained;
there are no infinite sums, projective identifications, matrix-equality
quotients or completely-positive-map equality quotients.
\end{definition}
\begin{scope}
The presentation specifies marked arithmetic circuits. It asserts neither
completeness of its equations nor faithfulness of a realization. Higher
Clifford levels label multiplication generators; a fixed higher level is
not declared to be closed under composition.
\end{scope}
\provenance{D1322,D1323}{theory/sidequests/frobenius-hierarchy/amplitude-category.md}{theory/checks/frobenius\_hierarchy\_check.py (A1--A4,A7)}
{theory/verdicts/frobenius-hierarchy-category-r1.md}

\begin{definition}[Hilbert interpretation of the arithmetic presentation]
Interpret $Q_E$ as $H_E$, $C_i$ as its actual support-code subspace with
orthonormal basis $|i(a)\rangle$, and words as ordered Hilbert tensors, with
$H(\mathbf1)=\mathbb C$. Interpret $w,j,v$ as $W,J,V$ above, and use
\[
 F_E|x\rangle=|E|^{-1/2}\sum_y\psi_E(-xy)|y\rangle,\quad
 U_E|x\rangle=|x^p\rangle,
\]
\[
 M_E^{(d)}|x_1,\ldots,x_d,z\rangle
   =|x_1,\ldots,x_d,z+x_1\cdots x_d\rangle
\]
for $f,u,m$. Interpret $e_E$ as $|0\rangle$, $t_i$ as the unitary
$|a\rangle\mapsto|i(a)\rangle$ with codomain the code, and $c_i$ as its
inclusion into $H_E$. Swaps are ordinary tensor swaps, dagger is adjoint,
and coefficients use the named complex embedding. Denote this candidate
interpretation by $H$.
\end{definition}
\begin{scope}
The negative Fourier kernel and fibre normalization are part of the data.
No basis or character compatibility beyond the named trace construction is
silently imposed.
\end{scope}
\provenance{D1302,D1304,D1306,D1308,D1324}{theory/sidequests/frobenius-hierarchy/amplitude-category.md}
{theory/checks/frobenius\_hierarchy\_check.py (A1--A4)}{theory/verdicts/frobenius-hierarchy-category-r1.md}

\begin{theorem}[A small arithmetic amplitude category]\statusProved
For every finite named field diagram and coefficient datum above, the
presentation is a small linear dagger symmetric monoidal category, and
$H$ is a strong monoidal dagger functor. If the prime field belongs to the
diagram and $B=(\beta_1,\ldots,\beta_r)$ is a named prime-field basis of $E$,
the coordinate circuit
\[
 a_B=\sum_{x=\sum_lx_l\beta_l\in E}
        \left(\bigotimes_l b_{\mathbb F_p,x_l}\right)b_{E,x}^\dagger
       :Q_E\longrightarrow Q_{\mathbb F_p}^{\otimes r}
\]
is a source unitary. Its interpretation factors the field Weyl operators
into prime-field Weyls when positions use $B$ and momenta its trace-dual
basis.
\end{theorem}
\begin{scope}
The coordinate identification depends on its named basis. The theorem
asserts no reconstruction of the field from the unmarked linear envelope,
no faithfulness, and no full Clifford or all-process completeness.
\end{scope}
\provenance{FRP-CAT,D1327}{theory/sidequests/frobenius-hierarchy/amplitude-category.md}{theory/checks/frobenius\_hierarchy\_check.py (A1--A4,A7)}
{theory/verdicts/frobenius-hierarchy-category-r1.md}

Every relation is checked against the independent formulas. For example,
trace adjunction gives $F_EJ_i=V_iF_K$ and the logical Weyl relation with
$T_i(b)$ as momentum. The basis-ket equations prove the coordinate circuit
unitary by cancelling off-diagonal terms in $a_B^\dagger a_B$ and
$a_Ba_B^\dagger$. These identities supply arithmetic content to a source
whose equality was fixed before interpretation.

\begin{definition}[Source-certified processes and their equality]
Fix a countable naming universe containing the field-element labels and
closed under finite ordered tuples. All external tag sets and hidden index
sets are finite subsets of this universe, with products formed by ordered
tuples in that universe. This naming convention is presentation data.
A process object is a finite nonempty tagged family $X=(X_a)_{a\in A}$ of
quantum words. An arrow $k:X\to Y$ consists of finite hidden index sets
$J_{b,a}$, possibly empty, and source amplitudes
$k_{b,a,j}:X_a\to Y_b$. For each $a$ there must be finitely many source
arrows $h_{a,l}:X_a\to Z_{a,l}$, to arbitrary quantum words, and a finite
source derivation of
\[
 1_{X_a}-\sum_{b,j}k_{b,a,j}^\dagger k_{b,a,j}
       =\sum_lh_{a,l}^\dagger h_{a,l}.
\]
The certificate is evidence, not additional arrow data. An arrow is
normalized when its displayed squared-amplitude sum equals $1_{X_a}$.
Equality is a bijection of each hidden index set carrying amplitudes to
equal source arrows. It does not allow general Kraus mixing, zero deletion,
external-tag relabeling or equality merely of realized CP maps. Distinct
external tags remain distinct even when their quantum words agree.

The identity has amplitude $1$ on each diagonal block and empty lists
elsewhere. For $k:X\to Y$ and $l:Y\to Z$, composition has hidden triples
$(b,j,t)$ on block $(c,a)$ and amplitude $l_{c,b,t}k_{b,a,j}$.
Tensor uses Cartesian products of external and hidden index sets and
amplitudes $k\otimes l$, in the stated factor order. The tensor unit is
the singleton empty word. Structural maps use canonical tag bijections
and word swaps. A deterministic tag map $g:A\to B$ has identity amplitude
in block $(g(a),a)$ and requires the same quantum word on each routed
block; this includes forgetting a classical tag, but cannot change its
quantum type. A coefficient-scalar process is a singleton coefficient
amplitude on the empty word with the same certificate; general endomorphisms
of that object may have closed-diagram amplitudes or several hidden indices.
\end{definition}
\begin{scope}
Composition sums over an intermediate classical tag. A history to be retained
must be part of the final external tag. No arbitrary complex-state or
arbitrary-CP-map generators are stipulated.
\end{scope}
\provenance{D1325}{theory/sidequests/frobenius-hierarchy/process-category.md}{theory/checks/frobenius\_hierarchy\_check.py (A7)}
{theory/verdicts/frobenius-hierarchy-category-r1.md}

\begin{definition}[Ordinary-trace process interpretation]
Assign the block algebra
\[
 B_X=\bigoplus_a\operatorname{End}_{\mathbb C}H(X_a),\qquad
 \operatorname{Tr}_X(\rho)=\sum_a\operatorname{Tr}(\rho_a).
\]
A density has $\operatorname{Tr}_X(\rho)=1$. The target arrows are CP maps
that do not increase this trace. Tensor is ordinary finite-dimensional
tensor, distributed over tag blocks. The candidate process functor is
\[
 R(k)(\rho)_b=\sum_{a,j}H(k_{b,a,j})\rho_aH(k_{b,a,j})^*.
\]
The Born weight of tag $b$ is the ordinary trace of its block. Conditioning
divides by this weight only when positive and is not a linear process arrow.
Retaining a history means placing it in an external tag; ordinary composition
sums over its intermediate tag, unless that value was also routed to output.
\end{definition}
\begin{scope}
These are ordinary trace-one Hilbert densities, distinct from uniform
classical-trace and Hecke coefficient-trace densities. The density convention
is not changed by using an arithmetic source.
\end{scope}
\provenance{D1326}{theory/sidequests/frobenius-hierarchy/process-category.md}{theory/checks/frobenius\_hierarchy\_check.py (A7)}
{theory/verdicts/frobenius-hierarchy-category-r1.md}

\begin{theorem}[Positive independent composition]\statusProved
The source-certified processes form a small symmetric monoidal category.
Its normalized arrows form a symmetric monoidal subcategory. The map $R$
is a strong monoidal functor to trace-nonincreasing CP maps and takes
normalized arrows to trace-preserving maps. It is not faithful.
\end{theorem}
\begin{scope}
No comparison with a collective Hecke assembly map or a changing-parameter
endpoint is asserted. External outcomes and source amplitudes remain more
detailed than their CP interpretation.
\end{scope}
\provenance{FRP-CP}{theory/sidequests/frobenius-hierarchy/process-category.md}{theory/checks/frobenius\_hierarchy\_check.py (A7)}
{theory/verdicts/frobenius-hierarchy-category-r1.md}

The source certificate proves positivity without deciding equality by
realization. If $S=\sum k^\dagger k$ and $T=\sum l^\dagger l$, the tensor
deficit splits as
\[
 1\otimes1-S\otimes T=(1-S)\otimes1+S\otimes(1-T),
\]
a sum of squares built from the two certificates. Sequential composition
has the corresponding deficit $1-S+\sum k^\dagger(1-T)k$.
After interpretation these sums give nonnegative trace losses at every
matrix amplification. The singleton amplitude processes $\{1\}$ and
$\{-1\}$ are distinct at source but induce the same scalar CP map, proving
the stated failure of faithfulness directly.

\begin{definition}[Retained decoding, preparation and register circuits]
For a source isometry $s:X\to Y$, put $p_s=ss^\dagger$ and $q_s=1-p_s$.
The retained decoder has source $(Y)$ and target
$(\mathrm{success}:X,\mathrm{failure}:Y)$, with amplitudes $s^\dagger$ and
$q_s$. Its success-only arrow $d_s:(Y)\to(X)$ has amplitude $s^\dagger$
and residual certificate $q_s$; encoding has amplitude $s$. Apply this to
$j_i,v_i,c_i$, their composites and independent tensors, without stipulating
a reset completion. For $a:Q_K^{\otimes n}\to Q_K^{\otimes n}$, its code
lift is $a^{[i]}=t_i^{\otimes n}a(t_i^\dagger)^{\otimes n}$, in particular
for Frobenius, Weyl and multiplication gates. The Fourier comparison after
decoding applies $f_K$ on success and $f_E$ on failure.

Zero and basis preparations have the single amplitudes $e_E$ and $b_{E,a}$.
Quantum discard has one external output and hidden amplitudes
$b_{E,a}^\dagger$, indexed by $a\in E$. Code discard first applies
$t_i^\dagger$; word discard tensors the atom discards; tagged-family discard
sums all input blocks into the singleton empty-word output. A normalized
source-generated state followed by a retained decoder is a preparation/test
protocol. For a named basis $B$ with the prime field in the diagram, the
coordinate circuit is the displayed $a_B$ above, retaining $B$ as its label.
No basis-independent identification is imposed. The realized span of basis
matrix units contains all $R_p$-valued matrices; the field structure is in
the marked source data and relations, not reconstructed from that span.
\end{definition}
\begin{scope}
Preparations, effects and discards are those supplied by these source
circuits and certificates. A retained failure branch keeps its stated
quantum output type. Conditioning is used only at positive success weight.
\end{scope}
\provenance{D1327}{theory/sidequests/frobenius-hierarchy/amplitude-category.md; theory/sidequests/frobenius-hierarchy/process-category.md}
{theory/checks/frobenius\_hierarchy\_check.py (A7)}{theory/verdicts/frobenius-hierarchy-category-r1.md}

\begin{theorem}[Coherent successful descent with complete outcome records]
\statusProved
For the inclusion, trace-fibre and code isometries, retained decoding is
normalized, and successful decoding is trace-nonincreasing with probability
$\operatorname{Tr}(H(p_s)\rho)$. Successful decodes compose along field towers
and tensor independently. Full decoders keep their specified histories.
Fourier conjugation intertwines both success and failure branches, and
code-lifted Frobenius and multiplication intertwine their physical gates
through the encoding.
\end{theorem}
\begin{scope}
The complete two-stage tower instrument is not identified with the
two-outcome decoder of the composite. The tensor of complete instruments
has four outcomes, including partial successes. No fixed higher hierarchy
level is declared a category.
\end{scope}
\provenance{FRP-DESCENT}{theory/sidequests/frobenius-hierarchy/amplitude-category.md; theory/sidequests/frobenius-hierarchy/process-category.md}
{theory/checks/frobenius\_hierarchy\_check.py (A3,A4,A7)}{theory/verdicts/frobenius-hierarchy-category-r1.md}

For isometries $s:X\to Y$ and $t:Y\to Z$, the three retained tower histories
have amplitudes $q_t$, $q_st^\dagger$, $s^\dagger t^\dagger$, with output
words $Z,Y,X$. Their effects sum to
$q_t+tq_st^\dagger+tp_st^\dagger=1$. The all-success amplitude is
$(ts)^\dagger$. For independent isometries, the four effects are
$p_s\otimes p_t$, $p_s\otimes q_t$, $q_s\otimes p_t$, $q_s\otimes q_t$.
These equations explain why no reset convention is needed for composition.

\subsection*{A support code and its Fourier-dual trace code}
Take $\mathbb F_4=\mathbb F_2[\alpha]/(\alpha^2+\alpha+1)$.
The inclusion code is spanned by $|0\rangle,|1\rangle$, while the relative
trace fibres are $\{0,1\}$ and $\{\alpha,\alpha+1\}$. Thus
\[
 J|0\rangle=|0\rangle,\quad J|1\rangle=|1\rangle,\qquad
 V|0\rangle=\frac{|0\rangle+|1\rangle}{\sqrt2},\quad
 V|1\rangle=\frac{|\alpha\rangle+|\alpha+1\rangle}{\sqrt2}.
\]
On $(|0\rangle+|\alpha\rangle)/\sqrt2$, inclusion decoding succeeds with
probability $1/2$, returning the conditional state $|0\rangle$. The failure
also has probability $1/2$, retaining the state $|\alpha\rangle$ in the
larger register. The maximally mixed input $I_4/4$ likewise succeeds with
probability $1/2$. The Fourier relation identifies the whole instrument:
\[
 V^*F_4=F_2J^*,\qquad (1-VV^*)F_4=F_4(1-JJ^*).
\]
\provenance{FRB-TRANSFER,FRP-DESCENT}{theory/sidequests/frobenius-hierarchy/process-category.md}
{theory/checks/frobenius\_hierarchy\_check.py (A3,A7)}{theory/verdicts/frobenius-hierarchy-category-r1.md}

Two independent inclusion decoders have output dimensions $4,8,8,16$ in
the order success--success, success--failure, failure--success,
failure--failure. On the entangled state
\[
 (|0,0\rangle+|\alpha,\alpha\rangle)/\sqrt2,
\]
the first and last outcomes
each have probability $1/2$, and both partial successes have probability
zero. Joint success is computed by $P_J\otimes P_J$; multiplying the
marginal success probabilities would incorrectly give $1/4$.
In the tower $\mathbb F_2\subset\mathbb F_4\subset\mathbb F_{16}$,
the three history outputs have dimensions $16,4,2$. Only the all-success
branch is the direct composite reduction to the prime field.
\provenance{FRP-DESCENT}{theory/sidequests/frobenius-hierarchy/process-category.md}
{theory/checks/frobenius\_hierarchy\_check.py (A7)}{theory/verdicts/frobenius-hierarchy-category-r1.md}

\subsection*{The next comparison at an incidence endpoint}
The proved arithmetic formulas do not yet define a specialization at one.
A concrete next proposal separates an incidence parameter $q$ from a fixed
cyclotomic phase level $N$. For extension-field isotropic contexts the
arithmetic point is $q=p^r,N=p$; prime-field contexts after restriction of
scalars have different rank and incidence parameter $p$. The proposed
comparison must carry the actual multiplication tables, inclusion and trace
matrices, code constraints and Frobenius processes to a coupled positive
context algebra. The type-C decomposable corner and shuffle correspondence
provide a composition test; their nonparabolic block does not supply a
generic Hecke tensor inclusion.

The reference trace matters. With phase level fixed, a normal-basis cyclic
Frobenius has phase trace $N^{1-r}$. Substituting $p=1$ into the different
physical arithmetic expression $p^{1-r}$ would force trace one, and a
unitary of trace one for a faithful positive normalized trace
$\tau$, with $\tau(1)=1$, equals the identity. A viable next construction must name its positive trace and
comparison functor, preserve an observed arithmetic process, and lift the
relation diagrams as well as their individual gates. The specific next
candidate and six witness tests are recorded in the accompanying endpoint
proposal; none is an admitted deformation theorem.
\provenance{D1250--D1255,D1001--D1007}{docs/research-plans/frobenius-coupled-limit.md; theory/sidequests/f1-limit/type-c.md}
{next-stage witness specification}{proposal only}
